% Generated draft: requires mathematical and source review.
\documentclass[11pt]{article}
\usepackage[margin=25mm]{geometry}
\usepackage{fontspec}
\setmainfont{Latin Modern Roman}
\usepackage{amsmath,amssymb,mathtools,unicode-math}
\setmathfont{Latin Modern Math}
\usepackage{enumitem,needspace,xurl,hyperref}
\setlength{\parindent}{0pt}
\setlength{\parskip}{5pt}
\setlength{\emergencystretch}{3em}
\newcommand{\N}{\mathbb{N}}
\newcommand{\Z}{\mathbb{Z}}
\newcommand{\Q}{\mathbb{Q}}
\newcommand{\R}{\mathbb{R}}
\newcommand{\C}{\mathbb{C}}
\newcommand{\F}{\mathbb{F}}
\newcommand{\A}{\mathbb{A}}
\newcommand{\PP}{\mathbb{P}}
\newcommand{\E}{\mathbb{E}}
\newcommand{\eps}{\varepsilon}
\newcommand{\dd}{\,\mathrm{d}}
\begin{document}
\section*{UnsolvedMath ID: 801056}
\textbf{Source problem code:} catalogue entry
\subsection{Definitions and mathematical statement}
The supplied record uses a class or parameter region $\mathcal E_0$, a parameter $b$, mild dissipativity, and infinite renormalizability without defining them. Their intended meanings must be recovered from the cited source before a complete formulation is possible. In particular, no specific H\'enon family, parameter domain, or definition of mild dissipativity is fixed by this record.

If $X$ denotes the intended ambient topological space containing $\mathcal E_0$, its boundary is
\[
\partial_X\mathcal E_0
=\overline{\mathcal E_0}^{\,X}\cap\overline{X\setminus\mathcal E_0}^{\,X}.
\]
The choice of $X$ and its topology is unresolved. The record poses the following two questions.
\begin{enumerate}
\item Describe $\partial_X\mathcal E_0$. In particular, is it a piecewise smooth arc? The meaning of arc, including its endpoint convention, and the smooth structure needed for piecewise smoothness are not specified.
\item For every infinitely renormalizable mildly dissipative map $f\in\mathcal E_0$, does its sequence of renormalizations converge?
\end{enumerate}
For the second question, let $\mathcal R_n(f)$ denote schematically the $n$th renormalization of $f$. Once the renormalization construction, its normalizations, and its target topological space $Y$ have been specified, the requested conclusion has the form
\[
\forall f\in\mathcal E_0,\qquad
\bigl(f\text{ is infinitely renormalizable and mildly dissipative}\bigr)
\Longrightarrow
\exists g\in Y:\ \mathcal R_n(f)\longrightarrow g\quad(n\to\infty).
\]
Here $Y$ and the convergence topology are unresolved placeholders, not additional universally quantified choices. The record does not specify renormalization combinatorics, domains, or a candidate limit.

The source motivates this question by proposing to implement the strategy attributed to reference [dCLM] for $b$ close to $0$. It supplies neither a quantitative range for $b$ nor an explicit statement that this proximity condition restricts the universal question above; that qualification remains unresolved.
\subsection{Short English statement}
The source asks whether the boundary of $\mathcal E_0$ is a piecewise smooth arc and whether every infinitely renormalizable mildly dissipative map in $\mathcal E_0$ has convergent renormalizations. It mentions the strategy of [dCLM] for $b$ close to $0$. Definitions of the family, $\mathcal E_0$, renormalization, and convergence are missing, so this is a qualified statement-only draft.
\subsection{Sources}
\begin{itemize}
\item Supplied local UnsolvedMath registry record: permanent ID $801056$, exact source problem code \texttt{catalogue entry}, title ``Henon Maps: Open Problems --- Question 8.'' The registry attributes the problem to Berger and Bedford, with proposed year 2023, and labels it open. This label is supplied metadata, not an independently verified mathematical or literature assessment.
\item Berger, Bedford, \emph{Henon Maps: A List of Open Problems} (2023), Question 8. Supplied source URL: \url{https://arxiv.org/abs/2312.03907}. The provenance records access on 2026-10-16 and extraction by \texttt{arxiv\_latex}; these are supplied provenance claims, not a claim of independent source access.
\item Reference [dCLM] is mentioned in the supplied statement as providing a strategy for $b$ close to $0$. No bibliographic details or separate link for [dCLM] are supplied.
\end{itemize}
\end{document}
